\documentclass[12pt]{article}

\usepackage{setspace}
\usepackage{amsmath,amssymb}
\usepackage{amsfonts}
\usepackage{graphicx}
\usepackage[pdftex,bookmarks=true,bookmarksopen=false,bookmarksnumbered=true,colorlinks=true,linkcolor=black,citecolor=black]{hyperref}
\usepackage[utf8]{inputenc}
\usepackage{float}
\usepackage{pdfpages}
\usepackage[top=3cm, bottom=2cm, left=3cm, right=2cm]{geometry}
\usepackage{indentfirst}	
\usepackage[brazil]{babel}
\renewcommand{\baselinestretch}{1.5}
\usepackage{subfigure}
\usepackage{multirow}
\usepackage{array}

%Estilo coluna centralizada e com largura variável
\newcolumntype{C}[1]{>{\centering\let\newline\\\arraybackslash}m{#1}}

%\usepackage{pstricks}%, egameps}
%\setlength{\textwidth}{17.2cm}
% \setlength{\textheight}{23cm}
%\addtolength{\oddsidemargin}{-22mm} 
%\addtolength{\topmargin}{-15mm} \addtolength{\evensidemargin}{-15mm}
%\setlength{\parskip}{1mm}
%\setlength{\baselineskip}{500mm}

\newtheorem{theorem}{Theorem}[section]
\newtheorem{assumption}{Assumption}
\newtheorem{acknowledgment}{Acknowledgment}
\newtheorem{algorithm}{Algorithm}
\newtheorem{axiom}{Axiom}
\newtheorem{case}{Case}
\newtheorem{claim}{Claim}
\newtheorem{conclusion}{Conclusion} 
\newtheorem{condition}{Condition}
\newtheorem{conjecture}{Conjecture}
\newtheorem{corollary}{Corollary}[section]
\newtheorem{criterion}{Criterion}
\newtheorem{defn}{Definition}[section]

\newtheorem{example}{Example}[section]
\newtheorem{exercise}{Exercise}
\newtheorem{lemma}{Lemma}[section]
\newtheorem{notation}{Notation}
\newtheorem{problem}{Problem}
\newtheorem{proposition}{Proposition}[section]
\newtheorem{remark}{Remark}
\newtheorem{solution}{Solution}
\newtheorem{summary}{Summary}
\newenvironment{proof}[1][Proof]{\textbf{#1.} }{\rule{0.5em}{0.5em}}

\begin{document}

\begin{titlepage}
\begin{figure}[!htb]
\includegraphics[width=15.92cm, height=4.03cm]{logo_PPGEE.png}
\label{logo}
\end{figure}
\begin{center}
\par
\vspace{100pt}
\textbf{\Large Nome do Aluno}\\
\vspace{102pt}
\textbf{\Large Título}\\
\end{center}

\par
\vfill
\begin{center}
{\singlespacing{\normalsize João Pessoa - PB}\\
{\normalsize Março de \the\year}}
\end{center}
\end{titlepage}

\thispagestyle{empty}

\newpage
\begin{figure}[!htb]
\includegraphics[width=15.92cm, height=4.03cm]{logo_PPGEE.png}
\label{logo2}
\end{figure}
\begin{center}
\vspace{30pt}
\textbf{\LARGE Nome do Aluno}

\par
\vspace{50pt}
\textbf{\Large Título}
\end{center}

\par
\vspace{30pt}
\hspace*{175pt}\parbox{7.6cm}{\singlespacing {\normalsize Exame de Qualificação de Mestrado submetida ao Programa de Pós-Graduação em
Engenharia Elétrica do Instituto Federal de Educação, Ciência e Tecnologia da Paraíba, como requisito necessário à obtenção do grau de Mestre em Ciências no Domínio da Engenharia Elétrica.}}

\par
\vspace{40pt}
\hspace*{125pt}\parbox{10.0cm}{{\normalsize Área de Concentração: Eletromagnetismo Avançado}}

\par
\vspace{30pt}
\hspace*{125pt}\parbox{10.0cm}{{\normalsize Orientador: Nome do Orientador, Prof. Dr.}\\
{\normalsize Co-Orientador: Nome do Coorientador, Prof. Dr.}}

\par
\vfill
\begin{center}
{\singlespacing{\normalsize João Pessoa - PB, Março de \the\year}}\\
{{\normalsize \copyright Nome do Aluno -  aluno.estudante@academico.ifpb.edu.com}}
\end{center}

% \thispagestyle{empty}
%  \newpage
%  \noindent{\textbf{\large Agradecimentos}}
%  \vspace{1pt}

% Agradecimentos...

% \thispagestyle{empty}
%  \newpage
%     \vspace*{\fill}
% 	\begin{flushright}
% 		\textit{``Nada lhe pertence mais que seus sonhos.''\\
% 		(Friedrich Nietzsche)}
% 	\end{flushright}

\thispagestyle{empty}
\newpage
\begin{center}
\textbf{\normalsize Resumo}
\end{center}
\vspace{1pt}
   
\noindent Escreva o seu resumo aqui.
\\  \\
Palavras-Chave: Modelo, latex, overleaf.

\thispagestyle{empty}

\newpage
\begin{center}
\textbf{\normalsize Abstract}
\end{center}
\vspace{1pt}

\noindent The english version of your "Resumo".
\\  \\
Keywords: Template, latex, overleaf.

\thispagestyle{empty}
\newpage
\listoffigures
\newpage
\listoftables
\newpage
\tableofcontents
\newpage

\include{cap1} 
%\include{cap2} 
%\include{cap3} 
%\include{cap4} 
%\include{cap5} 
%\include{cap6}

\newpage
\centering
\bibliographystyle{ieeetr}
\bibliography{refs} 



% \newpage
% \section{Apêndices}

% \subsection{xx}


% \begin{figure}[H]
% %left bottom right top
% %\includegraphics[trim=-10mm 55mm 00mm 15mm, scale=0.5]{gráfico.pdf}

% \begin{center}
% Figura 3: gráfico.
% \end{center}
% \end{figure}

% \begin{figure}[H]
% %left bottom right top
% %\includegraphics[trim=-10mm 55mm 00mm 15mm, scale=0.5]{gráfico.pdf}

% \caption{gráfico.}

% \end{figure}

% \newpage{}
% \subsection{Algoritmo}

% \subsubsection{yy}
% \begin{verbatim}
%   %Simulação;

%   for i = 1:n;
%     M(:,i) = rand(n,1);
%   end

% \end{verbatim}

\end{document}